\documentclass[11pt,a4paper]{article}
\usepackage[T1]{fontenc}
\usepackage[utf8]{inputenc}
\usepackage[french]{babel}
\usepackage{lmodern}
\usepackage{amsmath,amssymb}
\usepackage[margin=2.6cm]{geometry}
\usepackage{booktabs}
\usepackage{enumitem}
\usepackage{hyperref}
\hypersetup{colorlinks=true,linkcolor=blue!50!black}

\newcommand{\NW}{\mathrm{NW}}
\newcommand{\code}[1]{\texttt{#1}}

\title{Rapport d'implémentation du modèle M2\\
\large (workspace \code{m2\_fable})}
\author{Agent de réplication (workspace \code{m2\_fable})}
\date{3 juillet 2026}

\begin{document}
\maketitle

\begin{abstract}
Le modèle M2 a été implémenté conformément à la séquence à 7 phases du
document de conception (§2.2), avec les conventions C1--C10 de la
spécification technique. La suite de 31 tests passe intégralement. Deux
résultats méthodologiques importants sont apparus \emph{pendant les tests de
l'outillage statistique} : (1)~le critère « l'exponentielle gagne l'AIC sur
le corps » du protocole §6.1, tel qu'opérationnalisé dans le prototype, est
structurellement biaisé \emph{contre} l'exponentielle (sur un échantillon
exponentiel pur tronqué à 95\,\%, gamma gagne avec $\Delta$AIC $\approx 31$) ;
(2)~le test LR loi de puissance / log-normale du script historique
\code{tail\_test.py} — celui qui fonde l'affirmation « LR $+566$ à $+694$,
$p<10^{-3}$ » de [E7c] — conclut aussi « loi de puissance, $R \approx +575$,
$p \approx 0$ » sur des données \emph{log-normales pures}. Les deux outils ont
été corrigés (MLE tronqués). Par ailleurs, un premier run baseline révèle que
la règle M2 de transfert des créances au prorata rend le carnet de contrats
\emph{non stationnaire} (prolifération combinatoire), alors même que la
population et la richesse totale sont stationnaires.
\end{abstract}

\section{Objectifs}
Implémenter fidèlement M2 (sans casser l'existant, dans un workspace isolé),
le tester, et préparer l'infrastructure d'expériences du protocole §6.

\section{Méthode : architecture}

\begin{center}
\begin{tabular}{ll}
\toprule
Module & Rôle \\
\midrule
\code{src/m2/config.py} & dataclass gelée, variantes nommées, sérialisation \\
\code{src/m2/entities.py} & population en structure de tableaux, id jamais réutilisé \\
\code{src/m2/contracts.py} & carnet de prêts, agrégats créances/dettes (I2--I3) \\
\code{src/m2/market.py} & phase 6 : rounds, règle $w^*(\rho)$, taux géométrique \\
\code{src/m2/bankruptcy.py} & phase 7 : cascades, conventions C5--C7, variantes \\
\code{src/m2/simulation.py} & séquence des 7 phases, séries, snapshots \\
\code{src/m2/metrics.py} & E/S : \code{series.csv}, \code{snap\_*.npz}, JSON \\
\code{src/m2/analysis.py} & corps (MLE tronqué), CSN, LR corrigé, anti-cohorte \\
\code{src/m2/plots.py} & figures standard depuis les fichiers d'un run \\
\bottomrule
\end{tabular}
\end{center}

Expériences : \code{experiments/m2/run\_\{baseline,validation,grid,ablation\}.py}.
Tests : \code{tests/test\_\{core,step\_and\_bankruptcy,analysis\}.py}, exécutés
par \code{python3 tests/run\_all.py} (pytest absent du venv partagé ;
les tests sont pytest-compatibles).

\section{Fidélité à la spécification et écarts assumés}
Toutes les règles du §2.2 du rapport de conception sont implémentées telles
qu'écrites : naissances Poisson à bilan $(w_0, d_0)$, choc
$\eta \sim \mathcal N(-\sigma^2/2, \sigma^2)$ i.i.d. par entité, extraction
$\alpha\sqrt w$, service des intérêts avec défaut sur paiement partiel,
dépréciation du seul capital réel, marché à $\lfloor N/k\rfloor$ rounds avec
règle symétrique $w^*(\rho)$, faillites en cascade avec transfert au prorata.
Les points laissés ouverts par le rapport ont été tranchés par les
conventions C1--C10 (spécification, §3) ; aucune n'ajoute de paramètre libre.
Contrairement au prototype E7c, \emph{aucun} héritage du WIP
($\theta$, \code{offer\_frac}) n'est présent.

\section{Résultats}

\subsection{Tests}
31 tests, 0 échec : entités et $\NW$ ; conservation du principal (I4) et des
intérêts (I5) ; défaut sur paiement partiel ; faillite simple (résiduel au
prorata, extinction de $d_0$) ; transfert C5 (y compris part au
débiteur-créancier, créancier mort) ; variante \code{cancel} ; cascade à deux
niveaux ; invariants I1--I3 et I8 sur run court ; reproductibilité bit à bit
par seed (I10) ; cas limites $N \in \{0,1\}$, $N<k$, $w=0$, richesse égale,
défaillants exclus du marché ; récupération d'exposants et de familles sur
données synthétiques contrôlées.

\subsection{Découverte 1 : le critère AIC du corps était biaisé}
Le protocole §6.1 ajuste exponentielle/gamma/log-normale/Fisk sur les
95\,\% inférieurs. Or ajuster des familles \emph{non tronquées} sur un
échantillon \emph{tronqué} pénalise la vraie famille : sur $5000$ tirages
exponentiels purs ($\theta=5$), la méthode du prototype donne
gamma $<$ exponentielle avec $\Delta$AIC $= 31$. Autrement dit, le verdict
[E7c] « le corps n'est pas exponentiel, gamma gagne l'AIC » pouvait être en
partie un artefact de la méthode d'ajustement. Correction : MLE tronqué
(pdf renormalisée par $F(c)$ au seuil de troncature $c$), vérifié par test
($\Delta$AIC $\le 2$ sur exponentielle pure, la bonne famille gagne sur
log-normale pure).

\subsection{Découverte 2 : le test LR historique est biaisé pro-Pareto}
\label{sec:lr}
\code{tail\_test.py} (réutilisé par le prototype et cité par le protocole
§6.2) compare la Pareto MLE à une log-normale ajustée \emph{par les moments
de $\ln x$ de la queue}, sans renormalisation par la troncature en
$x_{\min}$. Sur $20\,000$ tirages \emph{log-normaux purs}
($\mu=0$, $\sigma=1$), il conclut « loi de puissance » avec
$R = +575$, $p \approx 0$ — du même ordre que les valeurs $+566$ à $+694$
rapportées par [E7c] comme preuve de queue de Pareto. La correction (MLE de
la log-normale \emph{tronquée à gauche} en $x_{\min}$) rend le test sain sur
données log-normales ; sa puissance sur données Pareto pures est en revanche
faible (une log-normale tronquée à grand $\sigma$ imite une Pareto), ce qui
est une limite connue de l'approche (Clauset et al.). Conséquence :
\textbf{la préférence « loi de puissance $\gg$ log-normale » revendiquée par
[E7c] n'est pas établie par les logs existants} ; la validation devra
s'appuyer sur la stabilité de l'exposant et les diagnostics mécanistiques
plus que sur ce LR.

\subsection{Découverte 3 : la règle de transfert au prorata fait proliférer
les contrats}
Premier run baseline (seed 0, $T=2000$, $\lambda=10$) : population
stationnaire ($\approx 1630$), richesse totale plate… mais contrats actifs
$1244 \to 3736 \to 43\,977$ entre $t=1000$, $1500$ et $2000$. Cause : à
chaque faillite, chaque contrat prêté de la faillie est scindé entre tous
ses créanciers (règle §2.2 phase 7) ; les grosses créancières accumulent des
milliers de micro-contrats qui sont re-scindés à leur mort — croissance
combinatoire. Le prototype E7c ne l'a jamais vue parce qu'il \emph{annule}
ces contrats au lieu de les transférer. C'est un vice de la spécification
M2 telle qu'écrite : le carnet n'est pas stationnaire, le coût par pas croît,
et la microstructure du réseau de crédit dérive. L'ablation
\code{fail\_lender\_loans=cancel} est donc élevée au rang de variante de
référence, et la comparaison des deux règles devient un résultat à part
entière (rapport de validation).

\section{Limites}
Le runner de tests est un substitut minimal de pytest (pas de fixtures ni de
paramétrage). Les assertions runtime coûteuses (I3 complet) ne sont
vérifiées qu'en fin de run et dans les tests, pas à chaque pas. Le test
mécanistique $T \approx B_0/A_0$ repose sur des incréments à 1 pas entre
snapshots denses $[1500, 1520]$ et exclut les entités mortes dans
l'intervalle (biais de survie près du plancher, discuté dans le rapport de
validation).

\section{Prochaines expériences}
Baseline 3 seeds ($T=2000$), ablations (\code{cancel}, \code{destroy},
\code{nocredit}, \code{arith}, $\lambda=30$), grille §6.5
($\sigma \times \varepsilon_0/w_0 \times k$), balayage $k \in \{2,3,4,6\}$,
puis protocole §6 complet par fenêtres et par seed.

\section{Conclusion}
L'implémentation est conforme, testée et reproductible. Trois découvertes
faites \emph{avant} toute validation montrent la valeur de la position
sceptique : deux outils statistiques hérités biaisaient les conclusions dans
le sens de la thèse (corps « non exponentiel », queue « Pareto prouvée »),
et une règle de la spécification produit une non-stationnarité structurelle
que le prototype masquait. La validation qui suit utilisera les outils
corrigés et comparera systématiquement les deux règles de faillite.

\end{document}
